% ISMAR 2026 Paper — Hybrid Smart Glasses–Smartphone Indoor Navigation
% Using IEEE VGTC conference style (vgtc.cls)

\documentclass[review]{vgtc}           % review mode (double-blind)
%\documentclass{vgtc}                  % final (camera-ready)

\graphicspath{{figures/}{images/}{./}}

\usepackage{times}
\renewcommand*\ttdefault{txtt}
\usepackage{booktabs}
\usepackage{mathptmx}
\usepackage{multirow}
\usepackage{xcolor}
\usepackage{balance}

\onlineid{0}
\vgtccategory{Research}
\vgtcinsertpkg

%% ============================================================
%% TITLE
%% ============================================================
\title{Second Screen, Second Opinion: Supporting Smart Glass Navigation with a Handheld Companion Display Under Uncertainty}

%% ============================================================
%% AUTHORS (anonymized for review)
%% ============================================================
% For review: leave anonymous
% For camera-ready: uncomment and fill in
%\author{Author One\thanks{e-mail: author1@university.edu}\\
%        \scriptsize University Name
%\and Author Two\thanks{e-mail: author2@university.edu}\\
%        \scriptsize University Name}

%% ============================================================
%% TEASER (optional — uncomment for camera-ready)
%% ============================================================
%\teaser{
%  \centering
%  \includegraphics[width=\linewidth, alt={System overview showing smart glasses AR arrows and smartphone information hub used simultaneously during indoor navigation.}]{teaser}
%  \caption{A participant uses XREAL Air2 Ultra smart glasses for AR arrow guidance while cross-referencing contextual information on the Beam Pro smartphone hub during an ambiguous intersection, illustrating the hybrid interaction paradigm evaluated in this study.}
%  \label{fig:teaser}
%}

%% ============================================================
%% ABSTRACT
%% ============================================================
\abstract{%
Optical see-through smart glasses can place AR navigation cues
directly in the user's view, but their limited field of view and
display resolution make it difficult to inspect contextual
information when guidance becomes ambiguous. We present a hybrid
indoor navigation approach in which smart glasses provide continuous
directional guidance (\emph{where to go}) while a handheld companion
display serves as an on-demand verification channel for contextual
information (\emph{what is there}). In a within-subjects study
($N{=}24$), participants completed mission-based wayfinding tasks on
a fixed indoor route under two conditions: Glass-Only and Hybrid.
Four uncertainty triggers were embedded along the route (tracking
degradation, information conflict, low-resolution guidance, and
guidance absence) to elicit cross-verification behaviour. We assessed
mission accuracy, completion time, confidence--accuracy calibration,
and a novel Cross-Verification Index~(CVI) that captures whether
participants consult the companion display more during uncertainty
than during routine navigation. Results showed that Hybrid
participants selectively increased companion-display usage during
uncertainty segments ($\text{CVI} > 1$) and exhibited stronger
confidence--accuracy alignment, while maintaining comparable mission
accuracy with a modest time trade-off. This work contributes a
dual-device navigation system implemented on commercial hardware, an
experimental protocol with operationalised uncertainty triggers, and
an analytic framework for studying complementary device roles in AR
wayfinding under uncertainty.
}

%% ============================================================
%% KEYWORDS
%% ============================================================
\keywords{Augmented reality, indoor navigation, hybrid interaction, smart glasses, device switching, confidence calibration, cross-device interaction.}

\nocopyrightspace  % remove for camera-ready

%%%%%%%%%%%%%%%%%%%%%%%%%%%%%%%%%%%%%%%%%%%%%%%%%%%%%%%%%%%%%%%%
%%%%%%%%%%%%%%%%%%%%%% START OF THE PAPER %%%%%%%%%%%%%%%%%%%%%%
%%%%%%%%%%%%%%%%%%%%%%%%%%%%%%%%%%%%%%%%%%%%%%%%%%%%%%%%%%%%%%%%

\begin{document}

\firstsection{Introduction}
\maketitle

%% ----- 1. INTRODUCTION -----

Augmented reality (AR) navigation systems can overlay directional cues
such as arrows or paths directly onto the physical world, reducing the
cognitive cost of translating between a map and the environment~\cite{Qiu:2025:UAR}.
Head-mounted displays (HMDs), in particular optical see-through smart
glasses, keep the user's hands free while providing an egocentric
perspective that aligns guidance with the walking direction~\cite{Rehman:2017:ARI}.
However, current consumer-grade smart glasses face well-documented
constraints: narrow field of view, limited display resolution, and
restricted input modalities that make detailed information
retrieval cumbersome~\cite{Li:2023:ERH, AReading:2025:RSU}.

Smartphones, conversely, excel at rich information display, precise
touch interaction, and high-resolution content rendering, but they
demand visual and manual attention that competes with environmental
awareness~\cite{Putra:2025:AAN, Sharin:2023:GCS}. Prior work has
compared HMDs and handheld devices as \emph{alternatives} for AR
navigation~\cite{Rehman:2017:ARI, Qiu:2025:UAR}, and a growing body
of research has explored hybrid HMD--smartphone interaction for
non-navigation tasks such as reading~\cite{AReading:2025:RSU,
Li:2023:ERH} and general-purpose AR~\cite{Zhu:2020:BIS,
Speicher:2018:XDA}. Yet no study to date has examined the
\emph{simultaneous, complementary} use of smart glasses and a
smartphone for indoor wayfinding, nor has any work investigated how
users switch between these devices under conditions of navigational
uncertainty.

This gap is consequential. Indoor environments frequently present
situations where the system's directional guidance alone is
insufficient: ambiguous intersections, tracking failures, conflicting
signage, or destinations that require attribute-level information (e.g.,
room capacity, equipment) to identify. In such moments, a second
device can serve as a \emph{cross-verification} channel---an on-demand
resource that either confirms or enriches the primary guidance. Whether
users actually adopt this strategy, and whether doing so improves
their trust calibration and task outcomes, remain open empirical
questions.

We address these questions through three contributions:

\begin{enumerate}
\item \textbf{A hybrid navigation system} built on XREAL Air2 Ultra
smart glasses and Beam Pro smartphone, implementing a dual-device
architecture where the glasses provide ambient AR arrows
(\emph{where}) and the phone offers a three-tab information hub
(\emph{what}: interactive map, contextual cards, mission reference).

\item \textbf{A controlled experiment} ($N{=}24$, within-subjects)
comparing Glass~Only versus Hybrid conditions on a fixed indoor route
with five mission-based wayfinding tasks and four embedded uncertainty
triggers.

\item \textbf{An analytical framework} centred on a novel
Cross-Verification Index (CVI) and confidence--accuracy calibration,
grounded in distributed cognition~\cite{Hutchins:1995:CIW} and trust
in automation~\cite{Lee:2004:TAS} theories.
\end{enumerate}


%% ----- 2. RELATED WORK -----

\section{Related Work}

\subsection{AR Indoor Navigation}

AR-based wayfinding has matured considerably over the past decade.
A recent systematic review by Qiu et al.~\cite{Qiu:2025:UAR}
analysed 88 studies and found that AR navigation improves wayfinding
performance (75\% of studies), reduces cognitive load (92\%), and
enhances spatial learning (85\%). Handheld AR systems such as
NavARNode~\cite{Putra:2025:AAN} and GoMap~\cite{Sharin:2023:GCS}
demonstrate accurate real-time guidance but require continuous
phone-holding, introducing fatigue and reducing environmental
awareness. Head-mounted approaches like NavMarkAR~\cite{Hsu:2023:NAL}
free the user's hands and improve immersion but provide limited
contextual detail. Rehman and Cao~\cite{Rehman:2017:ARI} directly
compared Google Glass and smartphone AR navigation, finding comparable
task times but lower workload with the HMD. Critically, all these
studies treat the two device form factors as \emph{alternatives}; none
examine their \emph{simultaneous, complementary} use.

\subsection{Hybrid HMD--Smartphone Interaction}

Several systems have explored bidirectional interactions between
HMDs and smartphones. BISHARE~\cite{Zhu:2020:BIS} proposed a design
space for smartphone--HMD information sharing, demonstrating benefits
of seamless cross-device data flow. AReading~\cite{AReading:2025:RSU}
quantified display switching costs in hybrid AR reading interfaces,
finding that high-resolution smartphone content offsets the overhead
of gaze shifts. Li et al.~\cite{Li:2023:ERH} showed that smartphones
as assistive displays reduce HMD reading workload. The XD-AR
framework~\cite{Speicher:2018:XDA} addressed cross-device AR
development challenges, while Lindlbauer et al.~\cite{Lindlbauer:2019:CAO}
demonstrated context-aware adaptation of MR interfaces. However, these
contributions focus on stationary tasks (reading, object manipulation)
and do not address the mobile, decision-laden context of wayfinding.

\subsection{Trust Calibration and Uncertainty in Navigation}

Trust calibration---the alignment between a user's confidence in a
system and the system's actual reliability---is a critical factor in
effective human--automation interaction~\cite{Lee:2004:TAS}. When
navigational uncertainty arises (e.g., tracking failures, ambiguous
guidance), miscalibrated trust can lead to over-reliance on faulty
cues or premature abandonment of helpful guidance~\cite{LeeMoray:1992:TCA}.
Colley et al.~\cite{Colley:2024:TCA} studied trust calibration in
AR decision support but not in a navigation context. The distributed
cognition framework~\cite{Hutchins:1995:CIW} suggests that
distributing information across multiple representations---here,
across two devices---can support more robust decision-making. We build
on these foundations to hypothesise that a hybrid device configuration
enables a \emph{cross-verification loop} that recalibrates trust when
the primary guidance source becomes unreliable.


%% ----- 3. SYSTEM DESIGN -----

\section{System Design}

\subsection{Hardware Platform}

The system runs on an XREAL Air2 Ultra optical see-through AR
glasses~\cite{XREAL:2024:A2U} tethered to an XREAL Beam Pro spatial
computing terminal~\cite{XREAL:2024:BPR}. The Air2 Ultra provides
6DoF tracking, hand tracking, and spatial anchoring. The Beam Pro
acts as both the computational host (Android) and the secondary
display surface with a touchscreen interface.

\subsection{Dual-Device Architecture}

The system implements a clear functional separation inspired by
distributed cognition theory~\cite{Hutchins:1995:CIW}:

\begin{itemize}
\item \textbf{Glasses (``WHERE'' channel):} AR directional arrows
rendered in egocentric view using head-locked placement. The arrow
colour, opacity, and angular spread are modulated by the active
uncertainty trigger. A minimal HUD shows the current mission and
waypoint progress.

\item \textbf{Smartphone (``WHAT'' channel):} A three-tab information
hub providing (1) an \emph{interactive floor-plan map} with POI
markers and route overlay, (2) \emph{contextual information cards}
(POI details, comparison matrices, contextual hints) that are
pushed at waypoints or pulled on demand, and (3) a \emph{mission
reference panel} for re-reading briefings and accessing
mission-type-specific hints.
\end{itemize}

\subsection{Condition Implementation}

In the \textbf{Glass~Only} condition, the smartphone hub is rendered
as a head-locked WorldSpace canvas on the glasses; users interact via
hand-ray pointing. In the \textbf{Hybrid} condition, the hub appears
on the Beam Pro touchscreen and activates via lift-to-wake. Both
conditions display the same information content, isolating the effect
of device form factor and interaction modality.

\subsection{Uncertainty Triggers}

Four trigger types are embedded at predetermined waypoints along the
route to systematically introduce navigational uncertainty:

\begin{itemize}
\item \textbf{T1 -- Tracking Degradation:} AR arrow exhibits
$\pm5^{\circ}$ jitter for 6\,s followed by a 3\,s blackout, then
auto-recovers.
\item \textbf{T2 -- Information Conflict:} The AR arrow direction
contradicts a physical directional sign (the arrow is correct).
\item \textbf{T3 -- Low-Resolution Guidance:} The arrow becomes a
$30^{\circ}$ semi-transparent fan, indicating directional uncertainty.
\item \textbf{T4 -- Guidance Absence:} The arrow is removed entirely;
only a ``near destination'' text hint remains.
\end{itemize}

Each participant encounters all four trigger types across two
sequential mission sets, counterbalanced for order.


%% ----- 4. EXPERIMENT DESIGN -----

\section{Experiment Design}

\subsection{Participants and Design}

We plan a within-subjects study with 24 participants
(target: university students, 18--35 years, unfamiliar with the
building). Each participant completes both conditions
(Glass~Only, Hybrid) in counterbalanced order, each with a
different mission set (Set~1, Set~2) on the same fixed indoor route
(Route~B, ${\sim}200$\,m, 9 waypoints, 6 turns). Four
counterbalancing groups (6 participants each) control for condition
and mission-set order effects.

\subsection{Mission-Based Tasks}

Each condition includes five missions executed in a fixed order
(A1$\to$B1$\to$A2$\to$B2$\to$C1):

\begin{itemize}
\item \textbf{Type~A -- Direction Verify} (2$\times$): Navigate to a
destination, then answer a 4-choice verification question about an
observable attribute (e.g., room number, signage content).

\item \textbf{Type~B -- Ambiguous Decision} (2$\times$): Navigate
through an uncertainty trigger zone, make a directional choice, and
verify correctness. These missions coincide with triggers T1--T4.

\item \textbf{Type~C -- Information Integration} (1$\times$): Visit
two locations, compare attribute information, and answer a comparative
judgement question (e.g., which room better suits a 10-person meeting).
\end{itemize}

\subsection{Measures}

\subsubsection{Primary Dependent Variables}

\begin{itemize}
\item \textbf{Mission accuracy:} Proportion of correct verification
answers (auto-scored, 4-choice).
\item \textbf{Completion time:} From \texttt{MISSION\_START} to
\texttt{VERIFICATION\_ANSWERED} (seconds).
\item \textbf{Confidence--accuracy calibration:} Per-participant
Pearson correlation between self-reported confidence (7-point Likert
at each waypoint) and mission correctness.
\item \textbf{Cross-Verification Index (CVI):}
$\text{CVI} = r_{\text{trigger}} \mathbin{/} r_{\text{baseline}}$,
where $r$ is the rate of information-hub access events per minute.
CVI~$> 1$ indicates elevated hub usage during trigger zones.
\end{itemize}

\subsubsection{Secondary Measures}

\begin{itemize}
\item NASA-TLX~\cite{Hart:1988:DOW} (per condition)
\item System Trust Scale~\cite{Jian:2000:FET} (7 items, per condition)
\item Device switching frequency and duration (from event logs)
\item Post-experiment preference and open-ended responses
\end{itemize}

\subsection{Hypotheses}

\begin{description}
\item[H1 (Trust Calibration):] The Hybrid condition produces a higher
confidence--accuracy calibration index than Glass~Only, indicating
that smartphone cross-referencing supports more accurate trust
assessment.

\item[H2 (Task Performance):] The Hybrid condition yields comparable
or higher mission accuracy, with a modest increase in completion time
attributable to device switching.

\item[H3 (Cross-Verification Behaviour):] In the Hybrid condition,
CVI is significantly greater than 1, indicating that participants
selectively increase information-hub usage during uncertainty triggers
relative to baseline navigation.
\end{description}

\subsection{Procedure}

The study follows a 50-minute protocol per participant:
(1)~consent and demographics (5\,min), (2)~Condition~1 with device
setup, route navigation with 5 missions, and NASA-TLX/trust survey
(15\,min), (3)~Condition~2 with the alternate mission set and surveys
(15\,min), (4)~comparative preference survey (10\,min), and
(5)~debriefing (5\,min). An experimenter follows 3\,m behind,
observing but not intervening unless a safety risk arises.

\subsection{Data Collection}

All interaction events are logged to a 15-column CSV at millisecond
resolution, capturing 35+ event types including
\texttt{WAYPOINT\_REACHED}, \texttt{BEAM\_SCREEN\_ON/OFF},
\texttt{TRIGGER\_ACTIVATED}, \texttt{CONFIDENCE\_RATED}, and
\texttt{VERIFICATION\_ANSWERED}. Head rotation is sampled
continuously. Glass-view video is captured for post-hoc analysis.


%% ----- 5. ANALYSIS PLAN -----

\section{Planned Analysis}

\subsection{Confirmatory Tests}

\begin{itemize}
\item \textbf{H1:} Paired $t$-test on calibration indices (Hybrid vs.\
Glass~Only). Supplementary: pre- vs.\ post-trigger confidence shifts.
\item \textbf{H2:} Paired $t$-test on mission accuracy and completion
time. Supplementary: $2 \times 3$ repeated-measures ANOVA
(condition $\times$ mission type).
\item \textbf{H3:} One-sample $t$-test of CVI against 1. Pearson
correlation between hub access frequency and mission accuracy.
\end{itemize}

\subsection{Exploratory Analyses}

We will examine (a)~mission-type-specific device switching patterns,
(b)~trigger-type-specific confidence recovery trajectories,
(c)~user clustering by information access strategy (map-centric vs.\
card-centric), and (d)~learning effects across conditions.


%% ----- 6. EXPECTED CONTRIBUTIONS -----

\section{Expected Contributions}

This work makes three contributions to the ISMAR community:

\begin{enumerate}
\item \textbf{Empirical evidence} on how hybrid smart
glasses--smartphone interaction affects trust calibration and
wayfinding performance under navigational uncertainty---a configuration
not previously studied.

\item \textbf{A novel metric (CVI)} for quantifying strategic
cross-device information seeking in trigger-elicited uncertainty,
applicable to other hybrid AR evaluation contexts.

\item \textbf{A replicable experimental framework} including
operationalised uncertainty triggers, mission-based task design, and
a fully implemented dual-device system on commercially available
hardware (XREAL Air2 Ultra + Beam Pro).
\end{enumerate}


%% ----- 7. LIMITATIONS -----

\section{Limitations and Future Work}

The study uses a single indoor environment, limiting
generalisability to other building types. The within-subjects design
with two mission sets, while controlling individual differences,
introduces potential learning effects that counterbalancing mitigates
but cannot eliminate. The four uncertainty triggers are designer-defined
and may not capture all naturalistic sources of navigational
uncertainty. Future work will extend the paradigm to multi-floor
environments, adaptive trigger generation based on real-time tracking
quality, and longitudinal studies of hybrid device adoption.


%% ----- 8. CONCLUSION -----

\section{Conclusion}

We present a hybrid smart glasses--smartphone navigation system and a
within-subjects experimental design to investigate whether distributing
navigational information across two complementary devices---ambient AR
arrows on the glasses and an on-demand information hub on the
phone---improves trust calibration and decision quality when wayfinding
becomes uncertain. Grounded in distributed cognition and trust-in-automation
theory, and operationalised through four uncertainty triggers and a
novel Cross-Verification Index, this work addresses an unexplored
intersection of hybrid device interaction, indoor AR navigation, and
trust-based behavioural analysis.


%% ----- SUPPLEMENTAL MATERIALS -----

\section*{Supplemental Materials}
\label{sec:supplemental_materials}

Upon acceptance, we will release the experimental protocol, event
logging schema, analysis scripts, and anonymised datasets on a
public repository under a CC BY 4.0 license. The Unity project
source code will be made available for replication purposes.


%% ----- ACKNOWLEDGMENTS (hidden in review) -----

\acknowledgments{
Removed for review.
}

%% ----- REFERENCES -----

\bibliographystyle{abbrv-doi-hyperref-narrow}
\bibliography{references}

\end{document}
